\subsection{Normalized secondary divisors and their complete unit-class coordinates}

\begin{statement}{A.211}{Proved}
{An actual lower capitulation map and three necessary unit coordinates}
Retain the fields, generators, and actual unit factorization of A.208.
Let \(a\) be an actual product of \(r_2,r_3,r_4\), put
\(v=\sigma^2(a)/\sigma(a)\), and choose an actual unit
\(b\in U_{K_3}\) with \(N_{K_3/K_2}b=v\).
The Hilbert--90 primitive and principalizer can be normalized,
without changing \(b\), so that
\[
\begin{gathered}
 b=\gamma(w)/w,\quad w=ux,\quad u\in U_\Omega,\quad x\in M^\times,\\
 N_sw=f\sigma(a),\quad f\in F^\times,\quad
 \operatorname{div}_M(x)=j_hD,\quad
 3D=\operatorname{div}_{E_2}(f),\quad \sigma D=D.
\end{gathered}
\tag{A.211.1}
\]
The full ordinary ideal class \([D]\) belongs to
\(\ker(j_h)\cap A_2^-=A_2^-\).
For the primes \(\mathfrak R_{q,1},\mathfrak R_{q,2}\)
of \(E_2\) whose lower residues have respectively
\(\alpha=\alpha_{0,q}\zeta,\alpha_{0,q}\zeta^2\), define
\(\Delta_q=[\mathfrak R_{q,1}]-[\mathfrak R_{q,2}]\).
These are the two A.173 basis classes with the specified order
of their conjugate primes.  Then
\[
 [D]=k_{11}\Delta_{11}+k_{59}\Delta_{59},\qquad
 \psi_2(v)=2(k_{11},k_{59}),\qquad
 k_q=v_{\mathfrak R_q}(x)\pmod3,
\tag{A.211.2}
\]
where \(\mathfrak R_q\) is the corresponding prime of \(M\).
This class is independent of all choices and extends continuously
to the completed relative-unit group:
\[
 \Psi_{\rm div}:\mathcal R_1\longrightarrow A_2^-,
 \qquad
 \Psi_{\rm div}(v)=
 2\psi_2(v)_{11}\Delta_{11}+2\psi_2(v)_{59}\Delta_{59}.
\tag{A.211.3}
\]

Let \(\overline U_M=U_M/U_M^3\), equivalently the same quotient
of its pro-three completion, and let \(\Lambda\) be the complete
six-place character map of A.208.  The following are equalities
of full subspaces through actual field restriction:
\[
\begin{aligned}
 \overline U_M^{\sigma=h=1,\,\tau=-1}
   &=\langle[z],[\epsilon],[u_{11}],[u_{59}]\rangle_{\F_3},\\
 \mathcal C:=
 \overline U_M^{\sigma=h=1,\,\tau=-1}\cap\ker\Lambda
   &=\langle[\epsilon],[u_{11}],[u_{59}]\rangle_{\F_3}.
\end{aligned}
\tag{A.211.4}
\]
Both lists are bases, and the lower radicals are
\[
 u_{11}=\frac{2+\zeta\alpha}{2+\zeta^2\alpha},\qquad
 u_{59}=
 \frac{2+2\zeta\alpha-\zeta^2\alpha^2}
      {2+2\zeta^2\alpha-\zeta\alpha^2}.
\tag{A.211.5}
\]
Their restricted field classes have actual global-unit
representatives in \(M\).  Every normalized choice in (A.211.1)
gives the actual unit \(c=x^3/f\), with
\[
 [c]=e[\epsilon]+k_{11}[u_{11}]+k_{59}[u_{59}]\in\mathcal C,
 \qquad N_su=\sigma(a)/c.
\tag{A.211.6}
\]
The last two coordinates are exactly those in (A.211.2).
The value selected by this actual \(\Omega\)-unit norm equation,
for each specified input \(v_2,v_3,v_4\), remains unevaluated.
\end{statement}

\begin{proof}
\emph{An integral alignment norm.}
Use \(K,r_0,f,w\) of (A.208.21)--(A.208.22).
The certified anti-unit norms give \(N_\sigma a=1\).
In \(\Z[\gamma]/(\gamma^9-1)\), put \(N_s=1+\gamma^3+\gamma^6\) and
\[
 L_j=P_4(\gamma^{4^{j-1}})-1-N_s,\qquad
 P_4(X)=1+X+X^2+X^3\quad(0\le j\le2),
\tag{A.211.7}
\]
where the exponent of four is read modulo three.  In ascending
powers of \(\gamma\), the literal integer coefficient vectors
\[
\begin{aligned}
 R_0&=(0,0,0,0,0,-1,0,-1,-1),\\
 R_1&=(0,-1,-1,0,0,-1,0,0,0),\\
 R_2&=(0,0,-1,0,-1,0,0,0,-1)
\end{aligned}
\tag{A.211.8}
\]
satisfy, by direct multiplication,
\[
 L_j=(\gamma^{4^j}-1)R_j,\qquad
 \sum_{j=0}^2R_j=-(\gamma+2\gamma^2)N_s.
\tag{A.211.9}
\]
We use multiplicative group-ring notation.  Since
\(N_sw=f\sigma(a)\) and \(h\) fixes \(\sigma(a)\),
\(N_h(f)=N_hN_s(w)/(\sigma a)^3\).
Thus the operator on \(w\) in \(N_h(r_0)\) is exactly
\[
 N_h\bigl(P_4(\gamma)h-1-N_s\bigr)
                 =\sum_{j=0}^2L_jh^j.
\tag{A.211.10}
\]
The actual relation \(h^j(b)=b\) gives
\((h^jw)^{\gamma^{4^j}-1}=b\).
Substitute (A.211.9), then use \(N_sb=\sigma^2(a)/\sigma(a)\):
\[
\begin{aligned}
 N_h(r_0)
 &=b^{-(\gamma+2\gamma^2)N_s}(\sigma a)^3\\
 &=a^{-(\sigma+2\sigma^2)(\sigma^2-\sigma)}(\sigma a)^3\\
 &=a^{N_\sigma-3\sigma}(\sigma a)^3=1.
\end{aligned}
\tag{A.211.11}
\]
This is the actual norm in \(F'/F\).  Hilbert--90 supplies
\(y\in F'^\times\) with \(r_0=h(y)/y\).  Replace
\[
 w\mapsto w/y,\qquad x\mapsto x/y,\qquad f\mapsto f/y^3.
\tag{A.211.12}
\]
Because \(\gamma\) fixes \(F'\), this preserves \(b\).
It preserves \(u=w/x\) and \(c=x^3/f\) literally.
The relation \(r_0^3=h(f)/f\) makes the new \(f\) lie in \(F\);
substitution into the formula for \(K\) makes the new \(r_0\)
equal to one.  This proves normalization with the original norm target.

\emph{The lower ideal and its actual class.}
Continue with the normalized notation.  The divisor
\(J=\operatorname{div}_M(x)\) satisfies
\(3J=\operatorname{div}_M(f)\) and is fixed by both \(h,\sigma\).
The only ramified prime of \(M/E_2\) is the unique three-prime.
Its coefficient in \(J\) is \(3v_{F,3}(f)\), since
\(M/F\) is totally ramified of degree nine there.
Every other invariant orbit descends without a divisibility
condition.  Therefore \(J=j_hD\) for an actual divisor of \(E_2\),
and injectivity of ideal extension gives
\[
 3D=\operatorname{div}_{E_2}(f),\qquad
 \sigma D=D,\qquad j_h[D]=0.
\tag{A.211.13}
\]
The A.208 identity is \(f\tau(f)=a_0^3\), initially with \(a_0\in F'\).
The cube class of \(f\tau(f)\) in \(F\) is \(\tau\)-fixed.
The actual restriction kernel
\[
 \ker(F^\times/F^{\times3}\to F'^\times/F'^{\times3})
                     =\langle[\zeta]\rangle
\tag{A.211.14}
\]
is \(\tau\)-minus.  Their intersection is zero, since two is
invertible modulo three.  Hence \(f\tau(f)\) is a cube in \(F\).
All its cube roots differ by \(\mu_3\subset F\), so the same \(a_0\)
lies in \(F\).  Thus the A.208 relation for \(J+\tau J\) descends to
\(D+\tau D=\operatorname{div}_{E_2}(a_0)\).
Together with \(3[D]=0\), this puts the full ordinary class in \(A_2^-\).

The primary class of a \(\sigma\)-invariant divisor is represented
by the seven ramified-prime coefficients, because
\(\operatorname{Cl}(F)\{3\}=0\).  The unique three-prime has zero
primary class by A.173.  In the relation for \(D+\tau D\), the
right side has the degree-three ramification factor at each changed
prime.  The \(\tau\)-fixed branch consequently has coefficient zero
modulo three; the two exchanged branches have opposite coefficients.
This gives \([D]=k_{11}\Delta_{11}+k_{59}\Delta_{59}\).
The coefficients agree with those of \(J\), since \(M/E_2\) is
unramified at these primes.  The actual pure-residue formula
(A.208.5) proves (A.211.2).
Its independence follows from the A.184 homomorphism \(\psi_2\)
and the independence of the two \(\Delta_q\) in A.173.
The products of \(v_2,v_3,v_4\) generate the full pro-three completion
by A.159; the finite-valued homomorphism therefore extends uniquely
and continuously as (A.211.3).

\emph{The complete necessary unit cube-class space.}
All expressions inside cubes in the following identities are
actual global units, by the invariant-divisor equations:
\[
\begin{gathered}
 \sigma(c)/c=(\sigma(x)/x)^3,\qquad
 h(c)/c=(h(x)/x)^3,\\
 c\tau(c)=(x\tau(x)/a_0)^3,\qquad N_su=\sigma(a)/c.
\end{gathered}
\tag{A.211.15}
\]
The last identity implies \(\Lambda(c)=\Lambda(\sigma a)=0\).
Indeed \(u\) is an actual unit norm preimage, while the residue
of the \(E_2\) unit \(\sigma a\) lies in \(\F_{q^2}\), on which
the \(\F_{q^6}\) cubic character of A.208 is zero.
This proves \([c]\in\mathcal C\).

Here is an explicit complete computation of that space.
Order the 27 basis columns as (A.199.8), and let \(R\) be their
invertible residue matrix at (A.199.7).  Evaluate those same units
using, respectively,
\[
 \sigma:(t,z)\mapsto(tz^3,z),\quad
 h:(t,z)\mapsto(t,z^4),\quad
 \tau:(t,z)\mapsto(t,z^8).
\tag{A.211.16}
\]
Each cubic residue is expressed using the original row generator
\(z^3\).  If \(R_a\) denotes the resulting matrix, \(R^{-1}R_a\)
is the actual action on the unit quotient.  This defines every
entry without an unspecified lift.  The certificate checks all
group relations, including \(\tau\sigma\tau=\sigma^{-1}\),
\(h\sigma=\sigma h\), and \(h\tau=\tau h\).

Let \(T\) be the stacked 81-by-27 matrix
\((\sigma-I,h-I,\tau+I)\), with entries represented by \(0,1,2\).
A rank-23 minor has determinant \(-38080\).
After appending the six tame rows, a rank-24 minor has determinant
\(246784\).  Both are prime to three.  The exact one-based row and
column lists are retained with the full matrices in the certificate
data.  These minors give invariant dimension at most four and
\(\dim\mathcal C\le3\).

The lower radicals in (A.211.5) have actual unit representatives
in \(M\).  Their numerator and denominator prime generators have
absolute rational norm \(q^2\) in \(F\), unique supports of residue
degree two, and the indicated opposite conjugate primes.  Therefore
\[
 \operatorname{div}_{E_2}(u_q)
                 =3(\mathfrak R_{q,2}-\mathfrak R_{q,1}).
\tag{A.211.17}
\]
In particular the raw class \(\Delta_q\) is killed by three.
A.206 makes its full extension to \(M\) principal.  If \(x_q\)
generates the divisor \(j_h(\mathfrak R_{q,1}-\mathfrak R_{q,2})\),
then \(u_qx_q^3\) is an actual global unit.  A different generator
changes this representative by a unit cube.  The unit \(\epsilon\)
already lies in \(F\).

Evaluate these two rational functions and the A.170 formula for
\(\epsilon\) at all 27 residue maps and apply \(R^{-1}\).
The resulting three coordinate columns have a determinant-eight
minor on rows \(2,1,3\).  Every column is killed by \(T\) and by
the six tame rows.  They therefore fill \(\mathcal C\).
The class \([z]\) is fixed by \(\sigma,h\) and is \(\tau\)-minus
modulo cubes, since \(h(z)/z=z^3\).  Its tame column is the nonzero
\(e_1\) column of (A.208.18).  These four classes fill the invariant
space, proving both equalities in (A.211.4).

\emph{The coordinate map back to the actual divisor.}
The extension \(M/F\) is the elementary Kummer extension generated
by \(z^3=\zeta\) and \(t^3=d^2\alpha\).  Thus
\[
 \ker(F^\times/F^{\times3}\to M^\times/M^{\times3})
              =\langle[\zeta],[d^2\alpha]\rangle.
\tag{A.211.18}
\]
The first line is \(\tau\)-minus, the second \(\tau\)-plus.
We proved that \([f]\) is minus in the lower field.  Since
\([c]=-[f]\) after restriction, (A.211.6) determines \([f]\)
as the negative of its three lower coordinates, up to \([\zeta]\).
At the chosen branch, \(u_q\) has valuation \(-1\);
\(\epsilon,\zeta\), and the other changed-prime radical are units.
Taking valuations therefore gives \(v_{F,q}(f)=k_q\pmod3\).
The relation \(x^3=fc\), with ramification degree three at \(q\),
gives \(v_{\mathfrak R_q}(x)=v_{F,q}(f)\).
This proves that the last two coordinates in (A.211.6) are
precisely the coefficients in (A.211.2).

The projection of \(\mathcal C\) onto its last two coordinates,
the class map \((k_{11},k_{59})\mapsto
k_{11}\Delta_{11}+k_{59}\Delta_{59}\), and multiplication by two
are therefore the explicit maps in the conclusion.
The actual equation \(N_su=\sigma(a)/c\), together with the actual
\(h\)- and \(\tau\)-fixedness of \(b\), retains the selection of a
class for an input.  The statement assigns no uncomputed value
to this selection or to its comparison with \(\Theta_h\eta\).

The executed certificate
\path{code/two_inert_secondary_invariant_units.py}
retains all action matrices and rank minors, the exact lower
radical columns, and the three integral identities (A.211.9).
Its complete inputs are the A.199 residue isomorphism and the
A.208 tame matrix, with their literal generator conventions.
\end{proof}
